\section{第1章 函数与极限}
\subsection{函数}
\begin{exercise}{1}{基础·定义域}
求 $f(x)=\arccos(2x+1)+\log_2(3x+2)$ 的定义域。
\end{exercise}
\begin{exercise}{2}{基础·复合函数}
设 $f(x)=\frac{x-1}{x+1}$。求 $f(-2)$ 和 $f(f(x))$，写出复合函数的定义域。
\end{exercise}
\begin{exercise}{3}{基础·换元求函数}
已知对一切 $x\ne0$，$f(1/x)=4/x+2x^2+1$。求 $f(t)$（$t\ne0$）及 $f(x^2+1)$。
\end{exercise}
\begin{exercise}{4}{基础·反三角函数}
设 $f(x)=\arcsin(4/x)$。求定义域及 $f(-4)$、$f(4\sqrt2)$。
\end{exercise}
\begin{exercise}{5}{常规·函数值域}
求 $f(x)=\sin\frac{\pi x}{2(1+x^2)}$ 的值域。提示：先求 $x/(1+x^2)$ 的取值范围。
\end{exercise}
\begin{exercise}{6}{基础·反三角函数值}
计算 $\arccos(\sqrt3/2)$、$\arctan(-1)$、$\operatorname{arccot}1$，其中 $\operatorname{arccot}$ 取值于 $(0,\pi)$。
\end{exercise}
\begin{exercise}{7}{常规·换元求函数}
已知 $f(x+1/x)=\frac{4x^2}{x^4+1}$（$x>0$）。令 $t=x+1/x$，求 $t$ 的取值范围及相应的 $f(t)$。
\end{exercise}
\begin{exercise}{8}{基础·分段函数}
设 $f(x)=\begin{cases}\frac{x}{\pi}|\sin x|,&|x|<\pi/3,\\0,&|x|\ge\pi/3.\end{cases}$ 求 $f(\pi/6)$、$f(-\pi/6)$、$f(2)$，判断其奇偶性。
\end{exercise}
\begin{exercise}{9}{常规·反解复合函数}
设 $\tan(2g(x))=x^2-1$ 且 $|g(x)|<\pi/4$。求 $g(x)$ 的表达式。
\end{exercise}
\begin{exercise}{10}{基础·分段表达式}
将 $\varphi(x)=\max\{x,x^2,1\}$ 写成分段函数，并求 $\varphi(-2)$、$\varphi(0)$。
\end{exercise}
\begin{exercise}{11}{基础·复合次序}
设 $f(x)=\begin{cases}-1,&x<0,\\1,&x\ge0,\end{cases}$、$g(x)=2x-1$。分别写出 $f(g(x))$、$g(f(x))$ 的分段表达式。
\end{exercise}
\begin{exercise}{12}{常规·分段复合函数}
设 $f(x)=2+|x|$、$g(x)=\begin{cases}x,&x<0,\\x^2,&x\ge0.\end{cases}$ 求 $f(g(x))$ 和 $g(f(x))$。
\end{exercise}
\begin{exercise}{13}{基础·分段函数运算}
设 $f(x)=\begin{cases}2x+1,&x<0,\\x^2+1,&x\ge0,\end{cases}$、$g(x)=x+1$。写出 $f+g$ 与 $fg$ 的分段表达式。
\end{exercise}
\begin{exercise}{14}{常规·反函数}
设 $f(x)=\frac{1-\sqrt{1+x}}{1+\sqrt{1+x}}$（$x\ge-1$）。求其反函数及反函数的定义域。
\end{exercise}
\begin{exercise}{15}{基础·奇偶性}
判断 $(x^2+1)\sin^3x$、$\ln(x+\sqrt{x^2+1})$、$\frac{2^{-x}-2^x}{2^{-x}+2^x}$ 的奇偶性。
\end{exercise}
\begin{exercise}{16}{常规·复合函数奇偶性}
设 $f,g:\mathbb R\to\mathbb R$，$f$ 为奇函数，$g$ 为偶函数。由定义判断 $f\circ f$、$f\circ g$ 的奇偶性。
\end{exercise}
\begin{exercise}{17}{常规·单调性的定义}
不用导数证明 $f(x)=x^3+2x+3$ 在 $\mathbb R$ 上严格递增。提示：对 $x_1<x_2$ 考察 $f(x_2)-f(x_1)$。
\end{exercise}
\begin{exercise}{18}{基础·有界性}
判断 $f(x)=1/(x-2)$ 在 $(2,3)$ 和 $(3,+\infty)$ 上是否有界。
\end{exercise}
\begin{exercise}{19}{基础·周期}
求 $f(x)=2+\cos^2(2x-1)$ 的最小正周期。
\end{exercise}
\begin{exercise}{20}{常规·取整函数}
记 $\lfloor x\rfloor$ 为不超过 $x$ 的最大整数。求 $f(x)=x-\lfloor x\rfloor$ 的值域及 $f(1.2)$、$f(-1.2)$，验证 $1$ 是它的周期。
\end{exercise}
\begin{exercise}{21}{基础·函数关系}
设 $f:(0,+\infty)\to\mathbb R$ 满足 $f(x)+2f(1/x)=x$。求 $f(x)$。提示：将 $x$ 换成 $1/x$ 后联立。
\end{exercise}
\begin{exercise}{22}{常规·函数综合}
设 $f(x)=(x-1)/(x+1)$。求定义域、零点和反函数，并写出反函数的定义域。
\end{exercise}
\subsection{数列的极限}
\begin{exercise}{23}{常规·极限定义}
用 $\varepsilon$--$N$ 定义证明 $\lim_{n\to\infty}n/(2n+1)=1/2$。提示：计算它与 $1/2$ 的差的绝对值。
\end{exercise}
\begin{exercise}{24}{基础·幂数列}
判断数列 $(1/2)^n$、$(-1/2)^n$、$(-1)^n$ 是否收敛，收敛时求极限。
\end{exercise}
\begin{exercise}{25}{基础·幂数列极限}
分别求 $\lim_{n\to\infty}n^{-2}$、$\lim_{n\to\infty}1/\sqrt n$、$\lim_{n\to\infty}n/(n+1)$。
\end{exercise}
\begin{exercise}{26}{基础·奇偶项极限}
设 $x_{2k-1}=-1/(2k)$、$x_{2k}=2^{-k}$（$k\ge1$）。分别求奇、偶子列的极限，并判断 $x_n$ 是否收敛。
\end{exercise}
\begin{exercise}{27}{基础·分段数列}
设 $x_n=(n+1)/(n+2)$（$n$ 为偶数），$x_n=1$（$n$ 为奇数）。求 $\lim_{n\to\infty}x_n$。
\end{exercise}
\begin{exercise}{28}{常规·子列判别}
设 $x_n\to-1$、$y_n\to1$，且 $z_{2n-1}=x_n$、$z_{2n}=y_n$。利用收敛数列的子列性质判断 $z_n$ 是否收敛。
\end{exercise}
\begin{exercise}{29}{基础·极限运算}
已知 $x_n\to2$。求 $|x_n|$、$x_{n+1}$、$1/x_n$ 的极限，其中倒数只考虑充分大的 $n$。
\end{exercise}
\begin{exercise}{30}{常规·振荡数列}
利用奇、偶子列判断 $x_n=(1+1/n)(-1)^n$ 是否收敛。
\end{exercise}
\begin{exercise}{31}{基础·绝对值极限}
设 $x_n\ge0$ 且 $x_n^2\to1$。求 $\lim_{n\to\infty}x_n$。提示：使用 $|x_n-1|=|x_n^2-1|/(x_n+1)$。
\end{exercise}
\begin{exercise}{32}{常规·数列极限综合}
设 $x_n=(2n+(-1)^n)/(n+1)$。求 $\lim_{n\to\infty}x_n$ 及 $\lim_{n\to\infty}(x_{n+1}-x_n)$。
\end{exercise}
\subsection{函数的极限}
\begin{exercise}{33}{常规·极限定义}
用 $\varepsilon$--$\delta$ 定义证明 $\lim_{x\to1/2}\frac{1-4x^2}{2x-1}=-2$。提示：先在 $x\ne1/2$ 时约分。
\end{exercise}
\begin{exercise}{34}{基础·无穷远极限}
分别求 $e^{2x-1}$ 与 $\arctan x$ 在 $x\to-\infty$、$x\to+\infty$ 时的极限。
\end{exercise}
\begin{exercise}{35}{基础·左右极限}
设 $f(x)=\begin{cases}(x+1)\sin x,&x<0,\\1-\cos x,&x\ge0.\end{cases}$ 求 $x=0$ 处的左右极限及双侧极限。
\end{exercise}
\begin{exercise}{36}{基础·分段极限}
设 $f(x)=\begin{cases}\ln(1-x),&x<1,\\1/\sqrt{2x+1},&x\ge1.\end{cases}$ 求 $x\to1^-$、$x\to1^+$、$x\to0$ 时的极限。
\end{exercise}
\begin{exercise}{37}{基础·分段函数极限}
设 $f(x)=\begin{cases}x+1,&x<0,\\3,&x=0,\\2-x,&x>0.\end{cases}$ 求 $f(0)$ 及 $x\to0$ 的左右极限，判断双侧极限是否存在。
\end{exercise}
\begin{exercise}{38}{常规·局部保号性}
设 $\lim_{x\to1}f(x)/x=-1$。先求 $\lim_{x\to1}f(x)$，再利用极限保号性说明 $f(x)$ 在 $1$ 的某去心邻域内为负。
\end{exercise}
\begin{exercise}{39}{常规·极限不等式}
设 $\lim_{x\to x_0}f(x)=A$、$\lim_{x\to x_0}g(x)=B$，且 $A<B$。证明在 $x_0$ 的某去心邻域内 $f(x)<g(x)$。提示：对 $g-f$ 使用保号性。
\end{exercise}
\begin{exercise}{40}{基础·振荡与夹逼}
求 $\lim_{x\to+\infty}\cos x/x$，并判断 $\lim_{x\to+\infty}\cos x$ 是否存在。
\end{exercise}
\begin{exercise}{41}{基础·局部有界性}
已知 $\lim_{x\to0}f(x)=2$。在极限定义中取 $\varepsilon=1$，证明 $f$ 在 $0$ 的某去心邻域内有界。
\end{exercise}
\begin{exercise}{42}{基础·复合极限}
已知 $\lim_{u\to0}f(u)=2$。求 $\lim_{x\to0}f(3x)$ 和 $\lim_{x\to0}f(x^2)$，其中只考虑 $x\ne0$。
\end{exercise}
\begin{exercise}{43}{常规·函数极限综合}
设 $f(x)=x\sin(1/x)+2x$（$x\ne0$）。用夹逼定理求 $\lim_{x\to0}f(x)$，并求 $\lim_{x\to0}xf(x)$。
\end{exercise}
\subsection{无穷小与无穷大}
\begin{exercise}{44}{基础·有界因子与无穷小}
用夹逼定理证明 $\lim_{x\to0}2x\sin(1/x)=0$。
\end{exercise}
\begin{exercise}{45}{基础·无穷大判定}
判断 $x^2$、$1/x$、$\sin x$ 在 $x\to+\infty$ 时是否为无穷大或无穷小。
\end{exercise}
\begin{exercise}{46}{基础·单侧极限}
分别求 $e^{1/x}$ 在 $x\to0^-$、$x\to0^+$ 时的极限，并指出在哪一过程中它是无穷小。
\end{exercise}
\begin{exercise}{47}{基础·无穷小与无穷大}
设 $f(x)=2x/(x-2)$。分别讨论 $x\to0$、$x\to2^-$、$x\to2^+$ 时的变化。
\end{exercise}
\begin{exercise}{48}{常规·乘积极限}
已知 $x_ny_n\to0$ 且 $x_n\to2$。求 $\lim_{n\to\infty}y_n$，写出使用商的极限法则的依据。
\end{exercise}
\begin{exercise}{49}{基础·有界因子估计}
求 $\lim_{x\to0}\sin x(2+\cos(3/x))$ 和 $\lim_{x\to+\infty}\arctan x/x$。
\end{exercise}
\begin{exercise}{50}{基础·无穷小判定}
判断 $1/x$（$x\to+\infty$）、$x\sin(1/x)$（$x\to0$）、$\ln(1+x)/x$（$x\to0$）是否为无穷小。
\end{exercise}
\begin{exercise}{51}{常规·极限保序性}
设 $x_n\to0$、$y_n\to1$、$z_n\to+\infty$。说明充分大的 $n$ 满足 $x_n<y_n<z_n$。提示：分别与 $1/2$ 和 $2$ 比较。
\end{exercise}
\begin{exercise}{52}{基础·由组合极限求极限}
已知 $\lim_{x\to0}(1+2x+f(x)/x)=4$。求 $\lim_{x\to0}f(x)/x$ 及 $\lim_{x\to0}f(x)$。
\end{exercise}
\begin{exercise}{53}{基础·无穷小运算}
求 $\lim_{x\to0}x^2\cos(1/x)$ 和 $\lim_{x\to0}x^2(1+\sin(1/x))$。
\end{exercise}
\begin{exercise}{54}{常规·倒数关系}
求 $\lim_{x\to0}1/x^2$ 和 $\lim_{x\to+\infty}1/(x^2+1)$，说明两式中无穷大与倒数无穷小的关系。
\end{exercise}
\subsection{极限的运算法则}
\begin{exercise}{55}{基础·极限四则运算}
已知 $f(x)\to2$、$g(x)\to3$（$x\to0$）。求 $f+g$、$fg$、$f/g$ 的极限。
\end{exercise}
\begin{exercise}{56}{基础·有限乘积极限}
求 $\lim_{x\to+\infty}(1-1/x)(2+2/x)(3-3/x)$。
\end{exercise}
\begin{exercise}{57}{基础·直接代入}
求 $\lim_{x\to1/2}(2\sin x+2\pi/x)/\cos2x$。
\end{exercise}
\begin{exercise}{58}{基础·有理化}
求 $\lim_{x\to1}(1/(x+1)-3/(x^3+1))$ 与 $\lim_{n\to\infty}(\sqrt{n+2\sqrt n}-\sqrt{n-\sqrt n})$。
\end{exercise}
\begin{exercise}{59}{基础·消因子与有理化}
求 $\lim_{x\to2}\frac{x^2+x-6}{x^2-x-2}$ 与 $\lim_{x\to0}\frac{x^2+x}{\sqrt{1+x}-1}$。
\end{exercise}
\begin{exercise}{60}{基础·最高次幂法}
求 $\lim_{x\to+\infty}\frac{2x^2+x\cos x}{x^3-x\sin x+1}$。
\end{exercise}
\begin{exercise}{61}{基础·和式极限}
求 $\lim_{n\to\infty}\sum_{k=1}^n\frac{k}{n^2+1}$ 与 $\lim_{n\to\infty}\sum_{k=1}^n2^{-k}$。
\end{exercise}
\begin{exercise}{62}{常规·望远镜求和}
求 $\lim_{n\to\infty}\sum_{k=1}^n\ln\frac{k+1}{k+2}$。提示：先合并对数。
\end{exercise}
\begin{exercise}{63}{常规·乘积对数化}
求 $\lim_{n\to\infty}\frac1{n^2}\ln\prod_{k=1}^n2^k$。
\end{exercise}
\begin{exercise}{64}{常规·有限乘积}
利用 $(1-a)\prod_{k=0}^n(1+a^{2^k})=1-a^{2^{n+1}}$，求 $\lim_{n\to\infty}\prod_{k=0}^n(1+2^{-2^k})$。
\end{exercise}
\begin{exercise}{65}{常规·单侧极限}
分别求 $\frac{x^2-1}{x+1}e^{1/(x+1)}$ 在 $x\to-1^-$、$x\to-1^+$ 时的极限。
\end{exercise}
\begin{exercise}{66}{常规·无穷远根式差}
分别求 $\sqrt{x^2+x+3}-\sqrt{x^2-x+3}$ 在 $x\to+\infty$、$x\to-\infty$ 时的极限。提示：有理化时使用 $\sqrt{x^2}=|x|$。
\end{exercise}
\begin{exercise}{67}{基础·幂指函数连续性}
求 $\lim_{x\to e}(1+\sin2x)^{\ln\ln x}$。
\end{exercise}
\begin{exercise}{68}{基础·幂指极限}
求 $\lim_{n\to\infty}(1/2)^n$ 和 $\lim_{n\to\infty}(3/2)^n$。
\end{exercise}
\begin{exercise}{69}{常规·极限反求参数}
已知 $\lim_{x\to2}\frac{x^2+ax+b}{x^2-4}=3$，求 $a,b$。提示：先令分子在 $2$ 处为零，再约分。
\end{exercise}
\begin{exercise}{70}{常规·重要极限反求参数}
求常数 $b$，使 $\lim_{x\to0}\frac{\sin2x}{\ln(1+x)}(\cos x-b)=4$。
\end{exercise}
\begin{exercise}{71}{常规·根式极限}
求 $\lim_{x\to-\infty}(4x+\sqrt{16x^2-8x-1})$。提示：有理化。
\end{exercise}
\begin{exercise}{72}{常规·根式差极限}
求 $\lim_{x\to+\infty}(\sqrt{x^4+2x^2+3}-x^2)$。
\end{exercise}
\begin{exercise}{73}{基础·极限运算变式}
已知 $f(x)\to1$、$f(x)+g(x)\to3$（$x\to0$）。求 $g(x)$ 和 $f(x)g(x)$ 的极限。
\end{exercise}
\begin{exercise}{74}{常规·根式差综合}
求 $\lim_{x\to+\infty}(\sqrt{x^2+4x+1}-x)$，并求该式在 $x\to-\infty$ 时的极限。
\end{exercise}
\subsection{极限存在准则与两个重要极限}
\begin{exercise}{75}{常规·夹逼定理}
利用 $0\le1/n!\le1/n$（$n\ge1$），求 $\lim_{n\to\infty}\sin(1/n)/n!$。
\end{exercise}
\begin{exercise}{76}{常规·和式夹逼}
求 $\lim_{n\to\infty}\sum_{k=1}^n\frac{k}{n^2+2n+k}$。提示：将每个分母分别与 $n^2+2n+1$、$n^2+3n$ 比较。
\end{exercise}
\begin{exercise}{77}{常规·振荡差式夹逼}
求 $\lim_{x\to+\infty}(\cos\sqrt{x+1}-\cos\sqrt x)$。提示：使用 $|\cos u-\cos v|\le|u-v|$，再有理化。
\end{exercise}
\begin{exercise}{78}{常规·递推数列}
设 $x_1=3$、$x_{n+1}=(x_n+4/x_n)/2$。证明 $x_n\ge2$ 且单调递减，进而求极限。
\end{exercise}
\begin{exercise}{79}{常规·单调有界数列}
设 $x_1=1$、$x_{n+1}=\sqrt{2+x_n}$。证明 $1\le x_n<2$ 且单调递增，求其极限。
\end{exercise}
\begin{exercise}{80}{常规·数列夹逼}
已知 $a_n=\int_0^1x^n\sqrt{1-x^2}\,\dd x$。利用 $0\le\sqrt{1-x^2}\le1$ 证明 $a_n\to0$。
\end{exercise}
\begin{exercise}{81}{基础·重要极限}
分别求 $\sin2x/(5x)$ 在 $x\to0$ 与 $x\to+\infty$ 时的极限。
\end{exercise}
\begin{exercise}{82}{基础·三角极限}
求 $\lim_{x\to0}\tan3x/(4x)$ 与 $\lim_{x\to0}\sin5x/(2x)$。
\end{exercise}
\begin{exercise}{83}{基础·反三角极限}
求 $\lim_{x\to0}\arcsin2x/x$ 和 $\lim_{x\to0}\arctan x/(2x)$。
\end{exercise}
\begin{exercise}{84}{基础·三角和差极限}
分别求 $\frac{2x-\sin3x}{2x+\sin4x}$ 在 $x\to0$ 与 $x\to+\infty$ 时的极限。
\end{exercise}
\begin{exercise}{85}{基础·换元极限}
求 $\lim_{x\to3}\frac{\sin(x^2-2x-3)}{x-3}$。提示：将正弦中的多项式因式分解。
\end{exercise}
\begin{exercise}{86}{基础·余弦极限}
求 $\lim_{x\to0}(1-\cos2x)/x^2$。
\end{exercise}
\begin{exercise}{87}{基础·数列重要极限}
求 $\lim_{n\to\infty}n^2\sin(3/n^2)$。
\end{exercise}
\begin{exercise}{88}{基础·第二重要极限}
求 $\lim_{x\to+\infty}(1-2/x)^{2x+1}$。
\end{exercise}
\begin{exercise}{89}{基础·商底数幂指极限}
求 $\lim_{x\to+\infty}\left(\frac{x+1}{x-1}\right)^{2x}$。
\end{exercise}
\begin{exercise}{90}{常规·二次尺度幂指极限}
求 $\lim_{x\to+\infty}\left(\frac{2x^2-1}{2x^2+1}\right)^{x^2}$。提示：把底数写成 $1$ 加无穷小。
\end{exercise}
\begin{exercise}{91}{基础·对数等价量}
求 $\lim_{x\to0}\sin3x/\ln(1+2x)$。
\end{exercise}
\begin{exercise}{92}{常规·取对数求极限}
求 $\lim_{x\to0}(1+\tan2x)^{-3/x}$。提示：先求其对数的极限。
\end{exercise}
\begin{exercise}{93}{常规·幂指极限}
求 $\lim_{x\to0}(\cos x)^{1/x^2}$。可使用 $\ln(1+u)\sim u$ 和 $1-\cos x\sim x^2/2$。
\end{exercise}
\begin{exercise}{94}{常规·指数和的极限}
求 $\lim_{x\to0}\left((2^x+8^x)/2\right)^{1/x}$。提示：取对数后使用 $a^x-1\sim x\ln a$。
\end{exercise}
\begin{exercise}{95}{基础·对数与重要极限}
求 $\lim_{n\to\infty}\ln[(1+1/n)^n]$ 和 $\lim_{n\to\infty}(1+1/n)^n$。
\end{exercise}
\begin{exercise}{96}{常规·幂指极限反求参数}
求实数 $a$，使 $\lim_{x\to+\infty}(1+a/x)^x=e^2$。
\end{exercise}
\begin{exercise}{97}{基础·夹逼变式}
求 $\lim_{n\to\infty}(2+\sin n)/n$。
\end{exercise}
\begin{exercise}{98}{常规·重要极限综合}
求 $\lim_{x\to0}(1+2x+3x^2)^{1/x}$，并说明充分小的 $x$ 下底数为正。
\end{exercise}
\subsection{无穷小的比较}
\begin{exercise}{99}{基础·等价无穷小}
当 $x\to0^+$ 时，写出 $e^{\sqrt x}-1$、$\ln(1+\sqrt x)$、$1-\cos\sqrt x$ 的等价无穷小。
\end{exercise}
\begin{exercise}{100}{基础·无穷小阶数}
当 $x\to0$ 时，写出 $\arctan2x$、$\sqrt{1+x}-1$、$1-\cos2x$ 的等价无穷小，并比较阶数。
\end{exercise}
\begin{exercise}{101}{基础·无穷远等价量}
求 $\lim_{x\to+\infty}x\frac{2x-3}{x^2-x+1}$，据此写出原分式的等价式。
\end{exercise}
\begin{exercise}{102}{基础·指数等价量}
求 $\lim_{x\to0}(e^{x^2}-1)/x^2$。
\end{exercise}
\begin{exercise}{103}{常规·三阶等价量}
已知 $\tan x-x\sim x^3/3$（$x\to0$），求 $\lim_{x\to0}(\tan2x-2x)/x^3$。
\end{exercise}
\begin{exercise}{104}{基础·乘积等价量}
求 $\lim_{x\to0}\frac{(2^{x^2}-1)\arctan2x}{x^3}$。
\end{exercise}
\begin{exercise}{105}{基础·等价关系反求参数}
求常数 $a$，使 $1/(ax^2+1)\sim1/(2x^2+1)$（$x\to+\infty$）。
\end{exercise}
\begin{exercise}{106}{常规·无穷小阶数比较}
当 $x\to0$ 时，求 $(\sqrt{1+x^2}-1)\ln(1+x^2)$ 的等价无穷小，并判断它是否为 $x^2$ 的高阶无穷小。
\end{exercise}
\begin{exercise}{107}{常规·乘积中的等价代换}
求 $\lim_{x\to0}\frac{(1-\cos x)(e^x-1)}{x^2\sin x}$。
\end{exercise}
\begin{exercise}{108}{基础·等价关系反求系数}
求常数 $a$，使 $\sin(ax)\sim3x$（$x\to0$）。
\end{exercise}
\begin{exercise}{109}{常规·对数极限反推}
已知 $f(x)\to0$ 且 $\lim_{x\to0}\ln(1+f(x))/x^2=3$。求 $\lim_{x\to0}f(x)/x^2$。
\end{exercise}
\begin{exercise}{110}{基础·已知极限运算}
已知 $\lim_{x\to0}(\sin6x+xf(x))/x=0$。求 $\lim_{x\to0}f(x)$。
\end{exercise}
\begin{exercise}{111}{常规·振荡项与等价量}
求 $\lim_{x\to0}\frac{4\sin x+3x^2\cos(1/x^2)}{(1+\cos x)\ln(1+2x)}$。
\end{exercise}
\begin{exercise}{112}{基础·根式等价量}
求 $\lim_{x\to0}(\sqrt{1+\sin x}-1)/x$。
\end{exercise}
\begin{exercise}{113}{基础·指数复合等价量}
求 $\lim_{x\to0}(e^{1-\cos x}-1)/x^2$。
\end{exercise}
\begin{exercise}{114}{常规·对数差化简}
求 $\lim_{x\to0}\frac{\ln(x^2+e^{2x})-2x}{x^2}$。提示：先把分子写为 $\ln(1+x^2e^{-2x})$。
\end{exercise}
\begin{exercise}{115}{常规·无穷远换元}
求 $\lim_{x\to+\infty}x\sin\ln(1+4/x)$。
\end{exercise}
\begin{exercise}{116}{常规·积分无穷小}
已知 $\int_0^x(e^t-1)\,\dd t=e^x-1-x$。利用 $e^x=1+x+x^2/2+o(x^2)$，求该积分在 $x\to0$ 时的等价无穷小。
\end{exercise}
\begin{exercise}{117}{常规·相消后的极限}
求 $\lim_{x\to0}(1/\sin x-1/x)$。提示：通分，可使用 $x-\sin x\sim x^3/6$。
\end{exercise}
\begin{exercise}{118}{常规·积分比值极限}
求 $\lim_{x\to0}\frac{x^3}{\int_0^x\ln(1+t^2)\,\dd t}$。可使用积分中值定理及 $\ln(1+t^2)\sim t^2$ 的估计，或使用洛必达法则。
\end{exercise}
\begin{exercise}{119}{基础·等价代换变式}
求 $\lim_{x\to0}\frac{\sin2x\,\ln(1+3x)}{x^2}$。
\end{exercise}
\begin{exercise}{120}{常规·二阶等价量综合}
已知 $\ln(1+x)=x-x^2/2+o(x^2)$。求 $\lim_{x\to0}\frac{\ln(1+x)-x}{1-\cos x}$。
\end{exercise}
\subsection{函数的连续性与间断点}
\begin{exercise}{121}{基础·连续函数值}
已知 $f$ 在 $1$ 处连续，且 $\lim_{x\to1}(f(x)-4)/(x-1)=2$。求 $f(1)$。
\end{exercise}
\begin{exercise}{122}{常规·连续性的定义}
用 $\varepsilon$--$\delta$ 定义证明 $f(x)=2x+1$ 在 $x=1$ 处连续。
\end{exercise}
\begin{exercise}{123}{基础·连续补定义}
设 $f(x)=(e^{2x}-1)/\sin4x$（$x\ne0$），$f(0)=k$，在 $0$ 的充分小邻域内考虑。求使 $f$ 在 $0$ 处连续的 $k$。
\end{exercise}
\begin{exercise}{124}{基础·分段连续性}
设 $f(x)=\begin{cases}\sin2x/\tan x,&x>0,\\2,&x=0,\\\ln(1-2x),&x<0.\end{cases}$ 各支在 $0$ 的充分小邻域内考虑。求左右极限，判断 $f$ 在 $0$ 处是否连续。
\end{exercise}
\begin{exercise}{125}{常规·振荡因子的连续性}
设 $f(x)=\begin{cases}(x-1)^2\cos(1/(x-1)),&x<1,\\(x+1)\ln x,&x\ge1.\end{cases}$ 判断 $f$ 在 $1$ 处是否连续。
\end{exercise}
\begin{exercise}{126}{常规·根式连续延拓}
为 $f(x)=\frac{\sqrt{1+x^2}-\sqrt{1-x^2}}{x\arcsin x}$（$0<|x|\le1$）补定义 $f(0)$，使其在 $0$ 处连续。
\end{exercise}
\begin{exercise}{127}{常规·数列极限生成函数}
对 $x>0$，设 $f(x)=\lim_{n\to\infty}\frac1n\ln(e^{nx}+x^2)$。求 $f(x)$。提示：将 $e^{nx}$ 提出。
\end{exercise}
\begin{exercise}{128}{常规·连续函数值反推}
设 $f$ 在 $0$ 处连续，且 $\lim_{x\to0}\frac{\sin(xf(x))}{x}=2$。求 $f(0)$。提示：可使用 $\sin u=u+o(u)$。
\end{exercise}
\begin{exercise}{129}{基础·有理函数间断点}
求 $f(x)=(x^2-1)/(x^2-3x+2)$ 的间断点并分类，对可去间断点作连续补定义。
\end{exercise}
\begin{exercise}{130}{常规·指数函数间断点}
求 $f(x)=1/(e^x-1)$ 的间断点及该点的左右极限，并分类。
\end{exercise}
\begin{exercise}{131}{常规·连续性反求参数}
设 $f(x)=\begin{cases}x\sin(1/x),&x<0,\\b,&x=0,\\a+\ln(1+3x)/x,&x>0.\end{cases}$ 求使 $f$ 在 $0$ 处连续的 $a,b$。
\end{exercise}
\begin{exercise}{132}{常规·点态极限函数}
设 $f(x)=\lim_{n\to\infty}1/(1+x^{2n})$。求分段表达式，并判断 $x=\pm1$ 处的连续性。
\end{exercise}
\begin{exercise}{133}{常规·夹逼生成函数}
利用 $\max\{1,x\}\le(1+x^n)^{1/n}\le2^{1/n}\max\{1,x\}$（$x\ge0$），求 $f(x)=\lim_{n\to\infty}(1+x^n)^{1/n}$，判断其连续性。
\end{exercise}
\begin{exercise}{134}{常规·分段连续性参数}
设 $f(x)=\begin{cases}\ln(1+2x)/x,&x<0,\\a,&x=0,\\(e^{bx}-1)/x,&x>0.\end{cases}$ 在 $0$ 的充分小邻域内考虑。求使 $f$ 在 $0$ 处连续的 $a,b$。
\end{exercise}
\begin{exercise}{135}{基础·连续函数运算}
求 $f(x)=e^{x/\sin x}$ 在 $x=0$ 处的极限，并作连续补定义。
\end{exercise}
\begin{exercise}{136}{基础·有理无理分段函数}
设 $f(x)=x^2$（$x\in\mathbb Q$），$f(x)=0$（$x\notin\mathbb Q$）。利用 $0\le f(x)\le x^2$ 判断 $f$ 在 $0$ 处是否连续。
\end{exercise}
\begin{exercise}{137}{常规·振荡函数连续补定义}
设 $f(x)=x^2\sin(1/x)$（$x\ne0$），$f(0)=c$。求使 $f$ 在 $0$ 处连续的 $c$，并判断 $\sin(1/x)$ 在 $0$ 处能否连续补定义。
\end{exercise}
\subsection{连续函数的运算与初等函数的连续性}
\begin{exercise}{138}{基础·连续函数商}
求 $\lim_{x\to0}\lg(100+x^2)/(1+\arcsin2x)$。
\end{exercise}
\begin{exercise}{139}{常规·连续函数复合}
设 $f$ 在 $0$ 处连续且 $f(0)=1$。求 $\lim_{x\to0}e^{f(x)}$、$\lim_{x\to0}\ln(1+f(x)^2)$。
\end{exercise}
\begin{exercise}{140}{基础·反函数连续性}
设 $f(x)=x^3$，求其反函数，说明反函数在实轴上连续。
\end{exercise}
\begin{exercise}{141}{常规·连续运算综合}
设 $f$ 在 $[a,b]$ 上连续且 $f(x)\ge1$。说明 $|f|$、$1/f$ 在 $[a,b]$ 上连续，并证明 $1/f$ 有界。
\end{exercise}
\subsection{闭区间上连续函数的性质}
\begin{exercise}{142}{常规·零点存在性}
证明方程 $x^3-3x+1=0$ 在 $(0,1)$ 内至少有一个实根。提示：对 $F(x)=x^3-3x+1$ 使用零点定理。
\end{exercise}
\begin{exercise}{143}{常规·零点存在与唯一性}
证明方程 $x=\frac12\cos x$ 在 $[0,1]$ 上有唯一实根。提示：考察 $F(x)=x-\frac12\cos x$ 的端点符号及单调性。
\end{exercise}
\begin{exercise}{144}{常规·平移差与零点}
设 $f$ 在 $[0,1]$ 上连续且 $f(0)=f(1)$。证明存在 $\xi\in[0,1/2]$ 使 $f(\xi+1/2)=f(\xi)$。提示：考察 $F(x)=f(x+1/2)-f(x)$ 的端点值。
\end{exercise}
\begin{exercise}{145}{常规·连续函数有界性}
设 $f$ 在 $[0,+\infty)$ 上连续且 $\lim_{x\to+\infty}f(x)=1$。证明 $f$ 有界。提示：将定义域分成一个闭区间和一个尾部区间。
\end{exercise}
\begin{exercise}{146}{常规·介值定理}
设 $f$ 在 $[a,b]$ 上连续。证明存在 $\xi\in[a,b]$，使 $f(\xi)=(f(a)+f(b))/2$。
\end{exercise}
\begin{exercise}{147}{基础·最值定理}
求 $f(x)=x(2-x)$ 在 $[0,2]$ 上的最大值与最小值，并说明最值存在的依据。
\end{exercise}
\begin{exercise}{148}{常规·零点定理应用}
设 $f$ 在 $[0,1]$ 上连续，$f(0)=-1$、$f(1)=2$。证明存在 $\xi\in(0,1)$，使 $f(\xi)=\xi$。提示：考虑 $F(x)=f(x)-x$。
\end{exercise}
\begin{exercise}{149}{基础·零点定理变式}
证明方程 $e^x+x-3=0$ 在 $(0,1)$ 内至少有一个实根。
\end{exercise}
\begin{exercise}{150}{常规·闭区间性质综合}
设 $f$ 在 $[a,b]$ 上连续且处处为正。用最小值定理证明存在 $m>0$，使 $f(x)\ge m$，进而说明 $1/f$ 在 $[a,b]$ 上有界。
\end{exercise}
